\section{Introduction}\label{sec:intro}

Over the past decade, China has built one of the world's largest digital economies. Broadband networks, mobile payments, e-commerce platforms, and cloud-based public services have expanded from leading coastal cities into inland and rural areas. This expansion has changed how firms organize production, how households access finance and markets, and how governments deliver public services \citep{goldfarb2019}. Yet the gains have not been evenly distributed. Per capita income, public-service capacity, and digital infrastructure remain uneven across provinces, and the gap between leading and lagging regions is a central concern of China's regional policy \citep{zhangky2019,worldbank2009}.

Digital factor allocation (DFA)---the combined availability and application of digital infrastructure, digital industry, digital finance, and digital human capital across regions---is a potential lever for regional coordinated development (RCD). Digital infrastructure lowers the cost of market access and information diffusion; digital finance extends payment, credit, and insurance services to households and small firms that were previously excluded; digital industry creates new employment and productivity opportunities; and digital human capital determines whether these opportunities are absorbed locally \citep{czernich2011,hjort2019,jack2014,zhang2019}. In this sense, DFA can support economic coordination, social coordination, and urban-rural coordination, the three core dimensions of RCD.

The same forces can also widen regional gaps. Digital investment tends to cluster where complementary skills, finance, and infrastructure are already available \citep{forman2012,akerman2015}. Platform markets reward scale and concentration, and automation can displace routine labor in regions that are least able to adjust \citep{acemoglu2020,autor2018}. Rural and lagging regions may gain access to digital services but lack the capability to convert access into productivity and income gains \citep{vandijk2005,norris2001,warschauer2003}. Whether DFA promotes or undermines RCD is therefore an empirical question, and the answer may depend on the stage of economic development and the local absorptive capacity.

Existing research on digitalization and regional development can be grouped into three strands. The first strand focuses on measurement. Studies construct digital economy indices from internet penetration, related employment, digital output, mobile internet, and digital financial inclusion \citep{zhao2020,guo2020}. These indices are useful for ranking regions, but they often aggregate heterogeneous digital factors into a single score and do not distinguish infrastructure, industry, finance, and human capital. The second strand examines the effects of digitalization on growth, productivity, employment, and inclusion \citep{czernich2011,hjort2019,huang2019,zhang2019}. The evidence is broadly positive, but the estimated effects vary across settings and are often sensitive to the level of development. The third strand introduces nonlinearity. Threshold models show that digitalization can have stage-dependent effects, with the marginal effect changing once a critical level of digital development, market efficiency, or technological progress is reached \citep{qi2025}. These studies suggest that linear estimates may mask important heterogeneity.

Three gaps remain. First, most studies use a composite digital economy index and do not examine whether the effect of DFA depends on complementary economic conditions. A composite index can be dominated by leading provinces and can obscure the separate roles of infrastructure, industry, finance, and human capital. Second, RCD is often measured by a single indicator, such as GDP growth or the urban-rural income ratio, rather than by a multidimensional index that reflects economic, social, and urban-rural coordination. Third, the threshold evidence for China is still limited, and the mechanisms through which DFA affects RCD are not well identified.

To fill these gaps, this study constructs a multidimensional DFA index and a multidimensional RCD index for 30 Chinese provinces from 2013 to 2021. The DFA index covers digital infrastructure, digital industry, digital finance, and digital human capital. The RCD index covers economic coordination, social coordination, and urban-rural coordination. Both indices are constructed with the entropy weight method, which avoids subjective weighting and preserves the information content of the underlying indicators. The study then applies a Hansen panel threshold model to examine whether the effect of DFA on RCD changes once economic development or financial development crosses a critical level. The analysis also examines regional heterogeneity and three plausible channels: industrial upgrading, financial deepening, and technological investment.

The contributions of this study are fourfold. (1) \textbf{Measurement}: It develops a transparent DFA index that distinguishes four digital factor categories and an RCD index that covers three coordination dimensions, allowing the analysis to move beyond a single composite score. (2) \textbf{Nonlinear identification}: It uses a panel threshold model to test whether the DFA effect depends on economic development and financial development, and it reports bootstrap tests and likelihood-ratio confidence intervals for the threshold estimates. (3) \textbf{Regional and channel evidence}: It shows how the effect varies across regions and identifies industrial upgrading, financial deepening, and technological investment as plausible channels. (4) \textbf{Policy relevance}: It provides a differentiated policy framework that sequences digital factor allocation with economic development and complementary reforms.

The results show that DFA is positively associated with RCD. Under two-way fixed effects, a one-unit increase in the DFA index is associated with an increase of 0.241 points in the RCD index (p < 0.05). The association is nonlinear. When per capita GDP exceeds the estimated threshold of 10.589 (about CNY 39,700), the marginal association rises from 0.506 to 0.677. A financial development threshold indicates that the marginal association declines at very high levels of financial depth, consistent with diminishing returns to digital finance. The association is strongest in western provinces, and the channel analysis suggests that industrial upgrading, financial deepening, and technological investment are plausible mechanisms. These findings are consistent with an absorptive-capacity interpretation: digital factor allocation is most effective when local economic conditions allow digital inputs to be combined with complementary skills, markets, and institutions.

The remainder of this paper is organized as follows. Section~\ref{sec:lit} reviews the literature and develops the hypotheses. Section~\ref{sec:data} describes the data, indices, and empirical strategy. Section~\ref{sec:results} reports the results. Section~\ref{sec:discussion} discusses the findings, policy implications, and limitations. Section~\ref{sec:conclusion} concludes.
